\documentclass[10pt]{article}
\usepackage{icmlko}

\icmlmeta{IP 파운데이션 월드모델}{967}{네 배는 인공물이었다}{앞 노트의 초록에서 가장 큰 수를 정정한다 --- 자기 열의 앞항을 빼 보는 데 한 줄이 들었고, 그 한 줄에서 간판 하나가 무너졌다}

\begin{document}

\twocolumn[
\icmlheader

\begin{abstract}
앞 노트(966)는 \emph{「집단 관심의 기억은 90일보다 길다」}는 명제를 내고 초록에
가장 큰 수를 실었다 --- \textbf{아이돌}은 판이 이미 보는 90일 수준이 $0.104$ 인데
\textbf{들뜸}이 $0.4230$ 이라 \textbf{「네 배」}라고. \textbf{그 수를 정정한다.}
정정의 근거는 새 자료가 아니라 \textbf{한 줄짜리 대조}다 --- 들뜸은 정의상
$\text{앞항}-\text{뒷항}$ 인데, 966 은 판이 이미 보는 \texttt{wiki\_level} 만 통제하고
\textbf{자기 열의 앞항(최근 90일)은 통제하지 않았다}. 아이돌은 부착 139 중
\textbf{108(77.7\%)}에서 긴 띠에 자료가 \textbf{0행}이고, 그 행에서 들뜸은
$\log(1{+}0)$ 을 빼는 것이라 \textbf{최근 수준 그 자체}다
($\mathrm{corr}=0.796$). 앞항을 같이 빼고 두 띠 모두 실자료가 있는 행만 남기면
$n$ 은 $139\to30$, $\rho$ 는 $+0.4230 \to \mathbf{-0.0199}$ $(p=0.91)$ 로
\textbf{부호까지 뒤집힌다}. 두 번째 정정은 자 자체다 --- 966 의 순위 함수는 동률을
평균내지 않아 \textbf{동률 블록에서 행 번호가 그대로 순위}가 되는데, 이 저장소의 판은
행이 결과로 정렬돼 있어 $\rho(\text{행번호},y)$ 가 두 도메인에서 정확히
$\mathbf{-1.0000}$ 이다. \textbf{정답표가 열에 샜다}. 동률평균으로 바꾸면 만화가
$0.1473 \to 0.0193$ 이 된다. \textbf{그런데 명제 자체는 산다.} 셋을 다 고치고
\textbf{Holm}$(\alpha{=}.05)$ 을 \textbf{측정 전에} 걸어도 10도메인 중
\textbf{셋}이 통과하고(세계애니 $+0.1747$, 애니 $+0.1912$, 게임 $+0.2311$),
\textbf{전부 양수이며 유의한 음수는 0}이다. 그리고 기전이 일관된다 --- 수준과
앞항을 뺀 뒤 \textbf{긴 띠 자체가 결과를 음으로 예측한다}. 이것은 후퇴가 아니라
\textbf{명제가 처음으로 자기 인공물과 갈라진 자리}다.
\end{abstract}
]

\section{무엇을 정정하는가}

정정 대상은 세 문장이다. 셋 다 966 이 \emph{자기 산출물에서 정직하게} 읽은
수이지만, 셋 다 \textbf{자를 한 번 더 조이면 남지 않는다}.

\begin{enumerate}\itemsep2pt
\item 초록의 \textbf{「아이돌은 수준 $0.104$ 인데 들뜸 $0.4230$ --- 네 배다」}
      $\to$ \textbf{정정: 인공물이다.}
\item 본문의 \textbf{만화 $0.1473$}(바닥을 넘은 다섯 중 하나) $\to$
      \textbf{정정: 동률평균으로 재면 $0.0193$ 이고 바닥을 못 넘는다.}
\item 헤드라인 \textbf{「10도메인 중 5」} $\to$ \textbf{정정: 다중비교 보정을
      걸고 규격을 조이면 \textbf{3} 이다.}
\end{enumerate}

\textbf{셋을 하나로 묶는 것은 「분석 자유도」다.} 966 은 사후 강건성 검사를
\textbf{하나}(문서 나이) 했고 그것은 옳았다. 하지 않은 것은 \textbf{가장 싸고 가장
결정적인 대조} --- \emph{내 열의 앞항을 통제에 넣으면 어떻게 되나} --- 였다.

\section{병 하나 --- 자기 열의 앞항}

들뜸의 정의는
\[
\text{들뜸}=\underbrace{\log(1+\overline{v}_{[s-90,\,s-1]})}_{\text{앞항}}
          -\underbrace{\log(1+\overline{v}_{[s-455,\,s-91]})}_{\text{뒷항}} .
\]
명제가 말하려는 것은 \textbf{뒷항이 나르는 정보}다. 그런데 966 은 통제 집합에
\texttt{wiki\_level}(\emph{다른} 원천 · \emph{다른} 창)만 넣었다. \textbf{앞항은
그대로 열 안에 남아 있다.}

그것이 언제 문제가 되나 --- \textbf{뒷항이 상수일 때}다. 뒷항이 상수면 들뜸은
앞항의 단조 변환이고, 명제가 아니라 \textbf{최근 수준}을 재게 된다. 그리고 뒷항이
상수가 되는 자리가 실제로 있었다.

\begin{center}\small
\begin{tabular}{lrr}
\hline
966 의 「유효」 개체 & \multicolumn{2}{r}{3{,}562}\\
\quad 두 띠 \emph{모두} 자료 0행 & 910 & (25.5\%)\\
\quad 긴 띠만 0행 $\Rightarrow$ 들뜸 $\equiv$ 앞항 & 357 & \\
\quad 긴 띠 0행 합계 & 1{,}267 & (35.6\%)\\
\hline
\end{tabular}
\end{center}

\textbf{두 띠가 다 0행인 910 행에서 들뜸은 날조된 $0.0$ 이다} --- 「기억이 없다」가
아니라 \textbf{「자료가 없다」}인데 $0$ 으로 들어갔다. 이 저장소는 이미 같은 원칙을
적어 두었다: \emph{「원천 개시일 밖은 0이 아니라 결측」}. 966 은 그 원칙을 띠
\textbf{밖}에만 적용하고 띠 \textbf{안}에는 적용하지 않았다.

\claimx{966 초록의 「아이돌 --- 네 배다」는 인공물이다. 자기 열의 앞항을 같이
통제하고 두 띠 모두 실자료가 있는 행만 남기면 $\rho$ 는 $+0.4230$ 에서 $-0.0199$
$(p=0.91)$ 로 부호까지 뒤집힌다.}{아이돌에서 앞항을 통제한 뒤에도 $\rho$ 가 양수로
유의하면 이 정정은 깨진다. 실측: 부착 139 중 긴 띠 0행이 108(77.7\%)이고
$\mathrm{corr}(\text{들뜸},\text{앞항})=0.796$ 이며, 실자료 행만 남기면 $n=30$ 에서
$\rho=-0.0199$ 다.}

\section{병 둘 --- 동률을 안 다루는 자가 정답표를 읽는다}

두 번째 병은 자료가 아니라 \textbf{자}에 있었다. 966 의 순위 함수는
\texttt{argsort().argsort()} 다. 안정 정렬이므로 \textbf{값이 같은 블록 안에서는
입력 순서가 그대로 순위}가 된다.

보통은 무해하다. 이 저장소에서는 아니다 --- \textbf{판의 행은 결과로 정렬돼
있다}. 실측하면 $\rho(\text{행번호},\,y)$ 가 \textbf{만화 $-1.0000$ · 세계애니
$-1.0000$} 이다. 그리고 만화는 부착 264행 중 \textbf{218행(82.6\%)에서 들뜸이 정확히
$0.0$}(두 띠 다 0행)이다. \textbf{그 218행의 순위는 자료가 아니라 정답표였다.}

이 실험실의 스피어만 정본은 \textbf{동률평균}이다(판을 자에 맞춘 것이 이미 기록에
있다). 정본으로 다시 재면 만화는 $0.1473 \to \mathbf{0.0193}$ $(p=0.75)$ 이고,
순열 바닥을 넘은 도메인은 \textbf{5에서 4}로 준다.

\claimx{동률을 평균내지 않는 순위 자는 판의 행 순서(=정답표)를 열에 넣는다.}{행
번호와 결과의 상관이 0 근처이면 이 주장은 힘을 잃는다. 실측: 두 도메인에서 정확히
$-1.0000$ 이고, 결과를 섞으면 12도메인 최댓값이 $0.0501$ 로 떨어진다.}

\section{그래도 명제는 산다 --- 좁혀서}

셋을 다 고친 규격(\textbf{동률평균} $+$ \textbf{앞항 통제} $+$ \textbf{두 띠 모두
실자료})으로 다시 잰다. 다중비교 보정은 \textbf{측정 전에} 정했다 ---
\textbf{Holm--Bonferroni}, $\alpha=0.05$, 분모는 그 규격에서 실제로 잰 도메인 수,
$p$ 는 라벨 순열 2{,}000회의 양측 값.

\begin{center}\small
\begin{tabular}{lrrrc}
\hline
도메인 & $n$ & $\rho$ & $p$ & Holm\\
\hline
세계애니 & 1{,}075 & $+0.1747$ & $0.0005$ & 통과\\
애니     & 362     & $+0.1912$ & $0.0010$ & 통과\\
게임     & 145     & $+0.2311$ & $0.0060$ & 통과\\
만화     & 45      & $+0.0435$ & $0.79$   & ---\\
아이돌   & 30      & $-0.0199$ & $0.91$   & ---\\
\hline
\end{tabular}
\end{center}

\textbf{분모 사다리}: 판 도메인 12 $\to$ 잰 도메인 10 $\to$ Holm 통과 \textbf{3}.
통과한 셋은 \textbf{전부 양수}이고 \textbf{Holm 을 통과한 음수는 0}이다.

그리고 \textbf{기전이 일관된다}. 같은 규격에서 뒷항 자체를 결과에 대고 재면 ---
수준과 앞항을 뺀 뒤 --- 애니 $-0.2340$, 세계애니 $-0.2171$, 게임 $-0.1976$ 으로
\textbf{음수}다. 읽으면 이렇다: \emph{같은 최근 수준이면, 오랜 기저가 두꺼울수록
결과가 나쁘다.} 들뜸이 양수로 결과를 예측하는 것과 \textbf{같은 한 가지}를 다른
쪽에서 본 것이다.

\claimx{명제는 좁아진 채로 산다 --- 판이 원리상 못 보는 띠가 결과 정보를 담는다.
분석 자유도가 만든 것은 넓이(5$\to$3)와 간판(아이돌 · 만화)이지 명제가 아니다.}{규격
D 에서 Holm 통과 도메인이 0 이거나, Holm 을 통과한 음수 도메인이 하나라도 있으면
이 주장은 깨진다. 실측: 통과 3 · 전부 양수 · 통과한 음수 0.}

\section{정직한 한계}

\textbf{첫째, 세계애니의 양수는 조건부 효과다.} 생 상관은
$-0.0078$ 로 음수인데 두 통제를 뺀 뒤 $+0.1747$ 이 된다(앞항과의 상관이
$-0.4038$ 이다). 억제 효과가 실재하는 신호일 수도, 통제 열의 잘못된 규격일 수도
있다. \textbf{이 표본은 그 둘을 못 가른다.}

\textbf{둘째, 판은 여전히 이 열을 판정할 자격이 없다.} 유보 3{,}775 중 부착이
$812(21.51\%)$ 뿐이고, 이 사이클이 규격 D 로 더 좁혔다. \textbf{채택 문턱은
$0.00353$ 하나}이며 --- 앞 노트가 쓴 $0.01055$ 는 \textbf{이미 철회된 진단
수치}이지 채택 크기가 아니다(정정) --- 어느 쪽이든 이 부착률에서는 잰다는 말이
성립하지 않는다.

\textbf{셋째, 이 정정 자체가 사후 선택이 아님을 보이는 방법은 하나뿐이다} ---
규격 · 보정 · 판정 규칙을 \textbf{측정 전에 단독 커밋}하는 것. 그렇게 했다.

\end{document}
